\subsection{Methodology}

The ScrumBan adoption for this project has been key to its successful execution so far. By combining selected characteristics of both Scrum and Kanban, this approach is well-suited to the project's needs.

Moving away from traditional Scrum ceremonies to brief meetings every two days has been a beneficial decision for the project. This approach enables continuous project monitoring and discards traditional ceremonies that provide low-value in a single-developer context. At the same time, restricting meeting duration to 30 minutes helps focus on critical issues and prevents overruns.

Regarding the WIP limit, it is not a restrictive constraint that blocks development progress. Quite the opposite, it has been a useful approach. Limiting the number of active tasks reduced the development boundary, allowing a more dynamic project flow.

Thanks to this methodology, the project's actual execution is kindly adhering to the scheduled timeline. At this stage, most of the tasks due today have been completed. Hence, it is not required to reassess the selected methodology.
